% !PRINTED 499
\weberpage{499}
%<<CONT
cients of $\varphi'$ are symmetric functions of the roots of the equation
\[ \frac{f(x)}{x - \Theta} = 0, \]
which are themselves contained in $\Omega$.

From the definition it follows that the norm of a product is equal to the product of the norms of the factors, and thus that, if $\varphi, \psi$ are functions in $\Omega$,
\weq{7}{N(\varphi\psi) = N(\varphi)\,N(\psi)}
holds.

\begin{center}
\S.\,136.

\gesperrt{The functionals of an algebraic field $\Omega$ and the extended field $\overline{\Omega}$.}
\end{center}

The variables that are used in the theory of algebraic numbers do not have the significance of symbols for variable sequences of numbers, as one is accustomed to from the theory of functions; rather, they are merely symbols of calculation without any independent meaning. In the functions of these variables what really matters is only the systems of coefficients, and the variables are used only in order to apply the well-known and familiar rules of literal calculation to these systems of coefficients. This does not, of course, exclude that occasionally the numbers are also considered which one obtains when the variables are replaced by certain numbers, e.g. by rational numbers.

\gesperrt{Accordingly we introduce into our considerations both integral and fractional functions of arbitrarily many variables, whose coefficients are numbers of an algebraic field $\Omega$, and stipulate that these functions are computed with as literal calculation prescribes.}

Every such function $\omega$ can be represented as a quotient of two integral functions in $\Omega$,
\weq{1}{\omega = \frac{\varphi}{\psi},}
where $\psi$ must always be assumed different from zero. Two such functions are only equal to one another:
\[ \frac{\varphi}{\psi} = \frac{\varphi_1}{\psi_1}, \]
%CONT>>
